\documentclass[11pt,a4paper]{article}
\usepackage[T1]{fontenc}
\usepackage[utf8]{inputenc}
\usepackage{lmodern}
\usepackage[a4paper,margin=25mm,headheight=15pt]{geometry}
\usepackage{xcolor,graphicx,booktabs,tabularx,longtable,multirow,arydshln,amsmath,amssymb,enumitem,multicol,fancyhdr,float,hyperref,natbib,makeidx,pdflscape,caption,array}
\definecolor{Navy}{HTML}{17324D}
\definecolor{Teal}{HTML}{0F766E}
\definecolor{Muted}{HTML}{586A7A}
\hypersetup{colorlinks=true,linkcolor=Navy,citecolor=Teal,urlcolor=Teal}
\pagestyle{fancy}\fancyhf{}
\lhead{Technical Writing Portfolio}\rhead{LaTeX Report}\cfoot{\thepage}
\renewcommand{\headrulewidth}{0.4pt}
\setlength{\parskip}{5pt}\setlength{\parindent}{0pt}
\setlist{itemsep=2pt,topsep=3pt}
\makeindex
\title{A Practical Portfolio for Technical Writing}
\author{Student Name \\ Department of Computer Science and Engineering}
\date{November 2026}
\begin{document}
\maketitle
\begin{center}\color{Muted}\small A worked example of structured writing, mathematics, figures, tables and references\end{center}
\tableofcontents
\listoffigures
\listoftables
\newpage
\section{Introduction}\label{sec:introduction}
Technical writing\index{documentation} communicates information so that a specific reader can understand, verify and use it. A good technical document is accurate, organised, readable and reproducible. This report demonstrates techniques that help achieve those goals.
\subsection{Objectives}
\begin{enumerate}
\item Organise a document using sections and subsections.
\item Apply consistent page and text styling.
\item Typeset mathematics, tables and images.
\item Connect content using labels, citations and an index.
\end{enumerate}
\subsection{A simple document workflow}
\begin{figure}[H]\centering
\fbox{\parbox{0.88\linewidth}{\centering\textbf{Plan}\quad$\rightarrow$\quad\textbf{Draft}\quad$\rightarrow$\quad\textbf{Review}\quad$\rightarrow$\quad\textbf{Publish}}}
\caption{A repeatable technical-writing workflow.}\label{fig:workflow}
\end{figure}
\section{Page styling and content formatting}\label{sec:formatting}
\subsection{Paper size, margins and page layout}
The document class and \texttt{geometry} package control page size and margins. This report selects A4 paper and a 25 mm margin. The \texttt{fancyhdr} package places the title in the header and the page number in the footer. The \texttt{pdflscape} package provides the \texttt{landscape} environment.
\begin{landscape}
\begin{table}[H]\centering
\caption{Example of a landscape page layout.}
\begin{tabularx}{0.88\linewidth}{lXl}
\toprule Setting & Purpose & Example\\\midrule
Paper size & Physical page dimensions & A4\\
Margins & Usable writing area & 25 mm\\
Orientation & Page direction & Landscape\\
Columns & Vertical reading lanes & Two columns\\
\bottomrule
\end{tabularx}
\end{table}
\end{landscape}
\subsection{Text styles, size, alignment and colour}
Normal text is used for paragraphs. \textbf{Bold text} highlights important terms; \textit{italic text} adds emphasis; and \texttt{monospaced text} represents commands or filenames. A colour example is {\color{Teal}\textbf{technical writing}}.
\begin{center}{\Large A centred line}\\[4pt]{\small A smaller supporting line}\end{center}
\subsection{Lists}
An unordered list is useful when order does not matter:
\begin{itemize}\item Clear headings.\item Consistent terminology.\item Correctly labelled figures and tables.\end{itemize}
A numbered list is appropriate when steps must be followed in order:
\begin{enumerate}\item Plan the content.\item Draft and format the document.\item Review, correct and publish it.\end{enumerate}
\subsection{Multiple columns}
The \texttt{multicol} package can divide a section into columns. This is useful for compact reference material, but long equations and wide tables may be better kept in one column.
\begin{multicols}{2}
\textbf{Column one.} Technical documents should use precise terms, define abbreviations and keep each paragraph focused on one idea.
\columnbreak
\textbf{Column two.} Lists, short examples and concise notes work well in columns. Check that the reading order remains clear.
\end{multicols}
\subsection{Writing mathematics}
Inline mathematics: $a^2+b^2=c^2$. A displayed, numbered equation:
\begin{equation}\label{eq:mean}\bar{x}=\frac{1}{n}\sum_{i=1}^{n}x_i.\end{equation}
Equation~\ref{eq:mean} gives the arithmetic mean of $n$ observations.
\begin{align}(a+b)^2&=(a+b)(a+b)\\&=a^2+2ab+b^2.\end{align}
A matrix and a piecewise function:
\[
A=\begin{bmatrix}1&2\\3&4\end{bmatrix},\qquad
f(x)=\begin{cases}x^2,&x\geq0,\\-x,&x<0.\end{cases}
\]

\section{Tables and images}\label{sec:tables-images}
\subsection{A basic table}
Table~\ref{tab:tools} uses booktabs rules for a clean academic style.
\begin{table}[H]\centering
\caption{Example tools used in a documentation workflow.}\label{tab:tools}
\begin{tabular}{lll}\toprule
Tool & Main purpose & Source format\\\midrule
LaTeX & Typeset reports & \texttt{.tex}\\
Doxygen & Generate API documentation & C++ comments\\
Beamer & Create slide presentations & \texttt{.tex}\\
Gnuplot & Plot numerical data & \texttt{.gp}\\\bottomrule
\end{tabular}
\end{table}
\subsection{Merged cells and dashed borders}
The \texttt{multirow} and \texttt{multicolumn} commands merge cells. The \texttt{arydshln} package supports dashed rules.
\begin{table}[H]\centering
\caption{A compact table with merged cells and dashed rules.}\label{tab:merged}
\begin{tabular}{|l|c|c|}\hline
\multirow{2}{*}{\textbf{Area}} & \multicolumn{2}{c|}{\textbf{Evidence}}\\
\cdashline{2-3}
& Source file & Output\\\hline
Documentation & C++ comments & HTML pages\\\hline
Presentation & Beamer source & PDF slides\\\hline
\end{tabular}
\end{table}
\subsection{Tables spanning multiple pages}
Use \texttt{longtable} when rows may continue over a page boundary. Its header can repeat on later pages.
\begin{longtable}{p{0.22\linewidth}p{0.66\linewidth}}
\caption{Example long-table content.}\label{tab:long}\\
\toprule\textbf{Stage}&\textbf{Description}\\\midrule
\endfirsthead
\toprule\textbf{Stage}&\textbf{Description (continued)}\\\midrule
\endhead
Planning & Identify reader, purpose, scope, constraints and expected output.\\
Audience analysis & Record what the reader knows and what they need to do.\\
Outline & Arrange sections and decide which diagrams, tables or examples are necessary.\\
Drafting & Organise information and write a first version.\\
Evidence & Check facts, measurements and code examples against reliable sources.\\
Figures & Number diagrams, add captions and refer to them in text.\\
References & Cite external ideas, data and quotations consistently.\\
Review & Check correctness, clarity, consistency and accessibility.\\
Publishing & Export to the required format and preserve editable sources.\\
Maintenance & Update the document when the system or requirements change.\\
\bottomrule
\end{longtable}
\subsection{Adding, scaling and rotating images}
The \texttt{graphicx} package provides \texttt{width}, \texttt{scale} and \texttt{angle} options:
\begin{verbatim}
\includegraphics[width=0.65\linewidth]{diagram.png}
\includegraphics[scale=0.5]{diagram.png}
\includegraphics[width=0.5\linewidth,angle=90]{diagram.png}
\end{verbatim}
Keep proportions unless there is a clear reason to stretch an image. Add a caption and label.

\section{Referencing and indexing}\label{sec:references}
\subsection{Cross-references}
Place a \texttt{\textbackslash label} after a heading, caption or numbered equation, then refer to it using \texttt{\textbackslash ref}. For example, Section~\ref{sec:formatting}, Table~\ref{tab:tools} and Figure~\ref{fig:workflow}. The \texttt{hyperref} package makes many references clickable.
\subsection{Citations and bibliography}
A citation connects a statement to its source. Version control\index{version control} keeps a traceable record of changes. This report uses BibTeX with \texttt{natbib}; see \citep{lamport1994latex} for a classic introduction to LaTeX. Entries are stored in \texttt{references.bib}.
\subsection{Index}
An index provides an alphabetical route to important terms. Add entries with \texttt{\textbackslash index\{Doxygen\}}. The \texttt{makeidx} package and MakeIndex program generate the index printed at the end.

\section{Conclusion}\label{sec:conclusion}
A well-structured document is easier to read, review and maintain. LaTeX separates content from presentation, while labels, citations and indexes make longer documents easier to navigate.
\section*{Acknowledgement}
This is a learning template. Replace sample identity and content with your own work before submitting it.
\bibliographystyle{plainnat}
\bibliography{references}
\printindex
\end{document}
